An assistant that cannot tell an under-specified request from a well-specified one will answer
both, and one of those answers is a guess dressed as an answer. We ask whether a model can
recognise that it lacks the information to act safely, and choose among \emph{acting},
\emph{asking}, \emph{refusing}, and \emph{deferring} according to the ambiguity and the stakes.
We build \carb, an 840-item four-way benchmark whose gold labels are inherited from the source
annotations of \coconot, \clamber, and \inthree rather than invented, and run five elicitation
regimes over identical items on five models.
Three results follow. First, \emph{structure beats confidence}: surfacing a typed action ontology
reaches 0.60--0.72 four-way accuracy against 0.23 uniform-random, and beats an
optimally-routed verbalised confidence by 0.16--0.28 on every model (all $p_{\text{Holm}}<10^{-8}$);
the scalar separates act from not-act (AUROC 0.79--0.87) but is at chance at separating \emph{why}
not (\defer \versus \refuse AUROC 0.44--0.56). Second, the failure is specific: models judge
``is this safe?'' at 0.89 and ``can I do this?'' at 0.94, but ``was I told enough?'' at 0.57.
Third, announcing high stakes moves the decision criterion ($\Delta c = -0.44$, $p<0.001$) without
improving discrimination ($\Delta d' = -0.11$, $p = 0.53$), and degrades it outright under a plain
ask-affordance. LoRA fine-tuning an open 4B model reaches 0.794 in-distribution but 0.470 out-of-source, where a
frozen probe on the untrained model already reaches 0.779. Once errors are priced, every prompted
frontier policy is worth less than always asking once a wrong action costs 5--10$\times$ a needless
refusal.
